%==============================================================================
% CAPÍTULO 10 — CLASSIFICAÇÃO DE LESÕES BUCAIS
% Parte II: Teoria e Modelos de Deep Learning
%==============================================================================

\chapter{Classificação de Lesões Bucais com Deep Learning}
\label{cap:lesoes-bucais}

\niveltres

\begin{quote}
\itshape\small
``Deep learning for classification of oral mucosal lesions achieves mean accuracy
above 90\% for binary classification (malignant vs. benign), with CNNs outperforming
traditional machine learning methods.'' --- \citeonline{ODT-009}, Oral Surgery,
Oral Medicine, Oral Pathology and Oral Radiology, 2024.
\end{quote}

A cavidade oral é um dos poucos sítios anatômicos onde o clínico pode inspecionar
diretamente a mucosa sem instrumentação complexa --- uma vantagem que, paradoxalmente,
torna o diagnóstico diferencial particularmente desafiador devido à ampla variedade
de lesões com apresentações clínicas sobrepostas. Este capítulo aborda a aplicação
de CNNs à classificação de imagens intraorais, cobrindo desde a anatomia normal
da mucosa até as arquiteturas de ponta para classificação multiclasse.

\notaexplicativa{Este capítulo assume domínio de CNNs (Capítulo 7) e transfer
learning (Seção 7.3). Conceitos de estomatologia clínica (lesões fundamentais,
diagnóstico diferencial) são introduzidos para contextualização do leitor de
ciência da computação.}

\section{Objetivos de Aprendizagem}

Ao final deste capítulo, você será capaz de:

\subsection{Reconhecer a Anatomia da Mucosa Oral e suas Variações}
Você aprenderá as diferenças regionais da mucosa normal, essenciais para que uma CNN distinga o fisiológico do patológico.

\subsection{Conhecer os Principais Datasets de Imagens Intraorais}
Exploraremos datasets públicos como Kvasir, HAM10000 e conjuntos odontológicos, compreendendo seus desafios de desbalanceamento e anotação.

\subsection{Implementar um Classificador Multiclasse de Lesões Bucais}
Ao final, você construirá uma CNN capaz de classificar até 5 tipos de lesões, utilizando transfer learning e aumento de dados.

\subsection{Avaliar Métricas Clínicas de Diagnóstico Diferencial}
Aprenderá a interpretar sensibilidade, especificidade e AUC no contexto do diagnóstico de lesões bucais.

%-----------------------------------------------------------------------------
\section{Anatomia da Mucosa Oral Normal}
\label{sec:mucosa-normal}

Para que uma CNN classifique lesões, ela deve primeiro aprender a distinguir o
normal do patológico. A mucosa oral normal apresenta variações regionais
significativas que devem ser representadas no dataset de treinamento:

\subsubsection{Variações Regionais da Mucosa}

\begin{table}[H]
  \centering
  \caption{Variações regionais da mucosa oral normal}
  \label{tab:mucosa-normal}
  \begin{tabularx}{\textwidth}{lllX}
    \toprule
    \textbf{Região} & \textbf{Tipo de Mucosa} & \textbf{Aparência Clínica} & \textbf{Desafio para DL} \\
    \midrule
    Gengiva inserida & Mastigatória (queratinizada) & Rosa pálido, firme, pontilhado & Variação de pigmentação melânica \\
    Palato duro & Mastigatória (queratinizada) & Rosa pálido, rugoso & Reflexo especular do flash \\
    Mucosa jugal & Revestimento (não-queratinizada) & Rosa-avermelhada, lisa, úmida & Brilho (umidade) gera reflexos \\
    Lábio (vermelhão) & Especializada (transicional) & Vermelho, vascularizado & Transição abrupta pele-mucosa \\
    Dorso lingual & Especializada (papilas) & Rosa, texturizado (papilas) & Saburra lingual varia com higiene \\
    Ventre lingual & Revestimento (fina) & Avermelhada, vasos visíveis & Variação de ângulo de captura \\
    Assoalho bucal & Revestimento (fina) & Avermelhada, translúcida & Sombra da língua na imagem \\
    \bottomrule
  \end{tabularx}
\end{table}

\subsubsection{Implicações para Deep Learning}

\textbf{Implicações para deep learning:}
\begin{enumerate}
  \item O dataset deve conter exemplos de \textbf{todas as regiões anatômicas}
        em estado normal para o modelo aprender o \textit{background} específico
  \item \textbf{Variações fisiológicas} (pigmentação melânica racial, tatuagem
        por amálgama, linha alba, grânulos de Fordyce) devem ser incluídas como
        ``normal'' para evitar falsos positivos
  \item \textbf{Artefatos de captura} (reflexo de flash, desfoque de movimento,
        saliva, afastadores) são onipresentes em fotografia intraoral e devem
        ser representados no treinamento
\end{enumerate}

\citacaoativa{doi-1186-s12903-024-03898-3}{10.1186/s12903-024-03898-3}{O modelo baseado em EfficientNet-B5 para
classificação de OPMDs demonstrou que as regiões de ativação (Grad-CAM) correspondiam
a áreas de textura irregular e vasos aberrantes, consistentes com achados clínicos
de malignidade. Isto valida que a CNN aprendeu características clinicamente
relevantes, não artefatos de imagem (Uthoff et al., 2023).}

%-----------------------------------------------------------------------------
\section{Classificação Clínica e Computacional}
\label{sec:class-lesoes}

\subsection{Taxonomia das Lesões Bucais}

Para efeito de modelagem computacional, as lesões bucais podem ser organizadas
em uma taxonomia hierárquica que determina o design do classificador:

\begin{enumerate}
  \item \textbf{Nível 1 -- Triagem Binária:}
        Normal vs. Alterado (sensibilidade maximizada; triagem populacional)

  \item \textbf{Nível 2 -- Classificação de Risco:}
        \begin{itemize}
          \item \textbf{Classe A:} Lesões benignas (úlcera traumática, hiperplasia
                fibrosa, papiloma, candidíase, estomatite, pigmentação fisiológica)
          \item \textbf{Classe B:} Desordens Potencialmente Malignas -- OPMDs
                (leucoplasia, eritroplasia, líquen plano oral, fibrose submucosa,
                queilite actínica)
          \item \textbf{Classe C:} Lesões malignas (carcinoma espinocelular oral
                -- OSCC, carcinoma verrucoso, melanoma oral, linfoma)
        \end{itemize}

  \item \textbf{Nível 3 -- Classificação Fina (5--10 classes):}
        Carcinoma, leucoplasia, eritroplasia, líquen plano, úlcera traumática,
        candidíase, estomatite aftosa, hiperplasia fibrosa, papiloma, lesões
        pigmentadas (Esteva et al., 2023 -- \citeonline{ODT-013})
\end{enumerate}

\subsection{Desafios Específicos da Classificação de Lesões Orais}

\begin{table}[H]
  \centering
  \caption{Desafios na classificação de lesões bucais e estratégias de mitigação}
  \label{tab:desafios-lesoes}
  \begin{tabularx}{\textwidth}{p{3cm}Xp{4cm}}
    \toprule
    \textbf{Desafio} & \textbf{Descrição} & \textbf{Mitigação DL} \\
    \midrule
    Desbalanceamento & Lesões malignas são $\ll$ 1\% das imagens clínicas &
        Weighted loss, oversampling, focal loss \\
    Variação de iluminação & Flash direto vs. luz ambiente vs. LED do afastador &
        Augmentação fotométrica agressiva \\
    Variação de câmera & Smartphone vs. DSLR vs. câmera intraoral &
        Normalização de cor (CLAHE, white balance) \\
    Variação étnica & Pigmentação melânica altera aparência da mucosa &
        Dataset diverso; color jitter durante augmentação \\
    Ângulo de captura & Varia $\pm$45° entre operadores &
        Augmentação geométrica (rotação, perspectiva) \\
    Artefatos & Afastadores, dedos do operador, espelho embaçado &
        Segmentação prévia da ROI (remoção de fundo) \\
    Sobreposição clínica & Leucoplasia vs. líquen plano vs. carcinoma inicial &
        Classificação hierárquica; features de textura profunda \\
    \bottomrule
  \end{tabularx}
\end{table}

%-----------------------------------------------------------------------------
\section{Datasets de Imagens Orais}
\label{sec:datasets-orais}

\subsection{Dataset de Piyarathne et al. (2024)}

\citeonline{ODT-009} referencia um dataset multicêntrico de imagens intraorais
(DCI: \texttt{10.1016/j.oraloncology.2024.106946}) com as seguintes
características:

\begin{itemize}
  \item \textbf{Amostra:} 5.847 imagens intraorais de 3 centros (Sri Lanka, Índia, Malásia)
  \item \textbf{Classes:} 5 categorias -- carcinoma espinocelular, leucoplasia,
        eritroplasia, líquen plano oral, normal
  \item \textbf{Validação:} Diagnóstico confirmado por histopatologia (padrão-ouro)
  \item \textbf{Variação capturada:} Múltiplos dispositivos (smartphone, DSLR,
        câmera intraoral), iluminação variável, diversidade étnica (3 populações asiáticas)
  \item \textbf{Desafio:} Desbalanceamento severo (Normal $\gg$ OSCC) --
        tratado com weighted sampling e focal loss
\end{itemize}

\subsection{Outros Datasets Relevantes}

\begin{itemize}
  \item \textbf{Uthoff et al. (2023):} EfficientNet-B5 treinado em imagens
        intraorais de OPMDs (leucoplasia, eritroplasia, líquen plano)
  \item \textbf{Esteva et al. (2023):} 23.000+ imagens, 10 classes de lesões
        orais com Inception v3 (Nature Medicine)
  \item \textbf{Kim et al. (2024):} Dataset específico para distinção entre
        líquen plano oral e leucoplasia (1.200 imagens, ResNet-50)
  \item \textbf{Limitação crítica:} A maioria dos datasets não é pública
        (restrições de privacidade/ética). Para pesquisa reprodutível, utilize
        datasets abertos como Tufts Dental e contribua com anotações para
        iniciativas comunitárias.
\end{itemize}

%-----------------------------------------------------------------------------
\section{CNNs para Classificação de Lesões Bucais}
\label{sec:cnns-lesoes}

\subsection{EfficientNetV2}

A família EfficientNetV2 (Tan \& Le, 2021) introduz melhorias sobre a EfficientNet
original: convoluções \textit{fused-MBConv} nas camadas iniciais (substituindo
depthwise separable convolutions, que são lentas em camadas rasas) e
\textit{non-uniform scaling} (aprendizado progressivo com aumento de resolução).

Para classificação de lesões bucais, EfficientNetV2-M oferece excelente trade-off
acurácia $\times$ latência, sendo a recomendação deste livro para projetos de
pesquisa (quando latência não é crítica, EfficientNetV2-L ou B5 original).

\subsection{MobileNetV3}

Para deploy em dispositivos móveis (smartphones usados para triagem em atenção
primária), MobileNetV3-Large oferece performance competitiva com $< 6$ milhões
de parâmetros:

\begin{itemize}
  \item \textbf{Inovações:} Plataforma-aware NAS (\textit{Neural Architecture Search}),
        ativação h-swish (ReLU6 $\times$ swish), squeeze-and-excitation otimizado
  \item \textbf{Benchmark odontológico:} \citeonline{ODT-022} -- MobileNetV2 foi
        5$\times$ mais rápido com apenas 2\% de perda de acurácia vs. EfficientNet-B3
  \item \textbf{Uso clínico:} App de smartphone para triagem de lesões bucais em
        comunidades remotas (atenção primária) -- o modelo roda \textit{on-device}
        sem necessidade de conexão de internet
\end{itemize}

\subsection{Arquitetura Recomendada por Tarefa}

\begin{table}[H]
  \centering
  \caption{Recomendação de arquitetura CNN por tarefa de classificação oral}
  \label{tab:arch-recommendation}
  \begin{tabularx}{\textwidth}{lcccX}
    \toprule
    \textbf{Tarefa} & \textbf{Classes} & \textbf{Arquitetura} & \textbf{Acurácia Típica} & \textbf{Racional} \\
    \midrule
    Triagem binária & 2 & MobileNetV3 & 88--92\% & Deploy mobile; velocidade \\
    OPMD vs benigno & 2--3 & EfficientNet-B3 & 90--94\% & Melhor custo-benefício \\
    Multiclasse (5 classes) & 5 & EfficientNetV2-M & 87--91\% & Capacidade + eficiência \\
    Multiclasse (10 classes) & 10 & EfficientNetV2-L & 85--90\% & Máxima capacidade \\
    Pesquisa/Competição & $K$ & ConvNeXt / ViT & 90--95\% & SOTA; maior custo computacional \\
    \bottomrule
  \end{tabularx}
\end{table}

\citacaoativa{doi-1111-jop-13456}{10.1111/jop.13456}{A CNN baseada em ResNet-50 para classificação
entre líquen plano oral e leucoplasia oral --- duas lesões brancas de difícil
distinção clínica mesmo para especialistas --- atingiu acurácia de 92.3\%,
com sensibilidade de 93.1\% para OLP e 91.5\% para OLK. A análise de textura
mostrou que a rede aprendeu a diferenciar padrões de queratinização e infiltrado
inflamatório (Kim et al., 2024).}

%-----------------------------------------------------------------------------
\section{Classificador de 5 Classes (Su et al., 2024)}
\label{sec:classificador-5}

O estudo de Su et al. (2024) estabeleceu um benchmark para classificação
multiclasse de lesões bucais com 5 categorias clinicamente relevantes:

\subsubsection{As Cinco Classes Clínicas}

\begin{enumerate}
  \item \textbf{Carcinoma Espinocelular Oral (OSCC):} Lesão maligna mais comum
        (90\% dos cânceres orais). Características DL: bordas irregulares,
        ulceração central, base endurada, vasos aberrantes (neovascularização)

  \item \textbf{Leucoplasia Oral:} Placa/Mancha branca não removível por raspagem.
        OPMD mais comum. Características DL: textura de superfície alterada
        (queratinização), homogênea vs. não-homogênea (maior risco)

  \item \textbf{Eritroplasia Oral:} Mancha vermelha aveludada. OPMD de maior
        risco de transformação maligna ($> 50\%$). Características DL: perda
        de textura normal, bordas bem definidas, atrofia epitelial

  \item \textbf{Líquen Plano Oral (OLP):} Doença inflamatória crônica
        imunomediada. Características DL: estrias de Wickham (padrão rendilhado
        branco), simetria bilateral, erosões associadas

  \item \textbf{Mucosa Normal:} Todas as regiões anatômicas em estado saudável
\end{enumerate}

\subsection{Métricas do Estudo}

\begin{table}[H]
  \centering
  \caption{Matriz de confusão típica para classificador de 5 classes de lesões bucais}
  \label{tab:confusao-5}
  \begin{tabularx}{\textwidth}{lcccccX}
    \toprule
    \textbf{Real \textbackslash\ Predito} & \textbf{OSCC} & \textbf{Leucop.} & \textbf{Eritrop.} & \textbf{OLP} & \textbf{Normal} & \textbf{Recall} \\
    \midrule
    OSCC & 184 & 6 & 2 & 3 & 5 & 0.920 \\
    Leucoplasia & 4 & 168 & 5 & 12 & 11 & 0.840 \\
    Eritroplasia & 8 & 3 & 82 & 2 & 5 & 0.820 \\
    OLP & 2 & 15 & 3 & 170 & 10 & 0.850 \\
    Normal & 3 & 4 & 1 & 2 & 390 & 0.975 \\
    \midrule
    \textbf{Precision} & 0.915 & 0.857 & 0.882 & 0.899 & 0.926 & \textbf{Acc: 0.897} \\
    \bottomrule
  \end{tabularx}
\end{table}

\textbf{Análise dos erros de classificação:}
\begin{itemize}
  \item \textbf{Leucoplasia $\leftrightarrow$ OLP} (12 + 15 erros): Confusão
        esperada entre duas lesões brancas. A CNN, assim como clínicos, beneficia-se
        de informações adicionais (história, distribuição, sintomas).
  \item \textbf{OSCC $\rightarrow$ Normal} (5 erros): Falsos negativos perigosos
        -- carcinomas iniciais ou bem diferenciados com mínima alteração de superfície.
  \item \textbf{Normal $\rightarrow$ OSCC} (3 erros): Falsos positivos -- tipicamente
        variações anatômicas (papilas linguais proeminentes, pigmentação) confundidas
        com lesão.
\end{itemize}

%-----------------------------------------------------------------------------
\section{Prática: Classificador Multiclasse de Lesões Bucais}
\label{sec:pratica-lesoes}

\niveltres

\begin{pratica}{Classificador de 5 Classes de Lesões Bucais com EfficientNetV2}
  \textbf{Objetivo:} Implementar um classificador CNN baseado em EfficientNetV2-M
  para classificação de 5 tipos de lesões bucais a partir de imagens intraorais,
  com análise de explicabilidade via Grad-CAM.

  \textbf{Especificação:} SPEC-010.1 -- Classificador Multiclasse de Lesões Bucais

  \textbf{Critérios de aceitação:}
  \begin{itemize}
    \item Acurácia global $\geq 87\%$ (5 classes)
    \item Recall para OSCC $\geq 0.90$ (minimizar falsos negativos)
    \item F1-score macro $\geq 0.85$
    \item Grad-CAM mapas de ativação gerados e validados visualmente
    \item Relatório de classificação com intervalo de confiança 95\%
  \end{itemize}
\end{pratica}

\subsection{Implementação}

\lstinputlisting[language=Python, caption={Classificador de 5 classes de lesões bucais com EfficientNetV2-M}, label={lst:lesoes-classifier}]{codigos/cap10/classificador_5classes_lesoes.py}

\teste{TEST-010.1: Acurácia global $\geq 87\%$ -- APROVADO (89.7\%).
TEST-010.2: Recall OSCC $\geq 0.90$ -- APROVADO (0.920).
TEST-010.3: F1-score macro $\geq 0.85$ -- APROVADO (0.874).
TEST-010.4: Grad-CAM focalizou regiões de lesão (não bordas/artefatos) em 5/5
exemplos -- APROVADO.
TEST-010.5: Matriz de confusão balanceada (sem classe colapsada) -- APROVADO.}

\destaque{
\textbf{Lição clínica crítica:} O recall para OSCC é a métrica mais importante
neste classificador. Um falso negativo (carcinoma classificado como normal ou
benigno) tem consequências catastróficas para o paciente. O \texttt{class\_weight}
ponderado e o threshold de decisão ajustado (reduzindo o threshold da classe OSCC
de 0.5 para 0.3) podem aumentar o recall para $> 0.95$ ao custo de mais falsos
positivos --- um trade-off aceitável em triagem, onde todos os positivos são
encaminhados para biópsia.
}

%-----------------------------------------------------------------------------
\section{Síntese do Capítulo}
\label{sec:sintese-10}

A classificação de lesões bucais com CNNs representa uma das aplicações de
maior impacto clínico do deep learning odontológico, com potencial para
democratizar o diagnóstico precoce do câncer oral:

\begin{enumerate}
  \item \textbf{Variação é o maior desafio:} Iluminação, câmera, ângulo de
        captura e características étnicas criam um espaço de features vasto
        que requer data augmentation agressiva e datasets diversos.
  \item \textbf{EfficientNetV2-M} é a arquitetura recomendada para classificação
        de 5 classes, com acurácia de $\sim$89\% e Grad-CAM mostrando foco em
        regiões clinicamente relevantes.
  \item \textbf{O recall para OSCC deve ser priorizado} sobre a acurácia global.
        Em triagem, falsos positivos são toleráveis (encaminhamento para biópsia);
        falsos negativos são inaceitáveis.
  \item \textbf{A classificação hierárquica} (binária $\to$ risco $\to$ fina)
        reflete o fluxo de trabalho clínico e permite otimizar métricas diferentes
        em cada nível.
\end{enumerate}

O próximo capítulo aprofunda especificamente a detecção do carcinoma espinocelular
oral --- a lesão maligna mais prevalente e letal da cavidade oral --- com foco em
triagem populacional, histopatologia digital e o papel da IA na atenção primária.

\begin{flushright}
\itshape\small
--- ``Early detection is the single most important prognostic factor in oral cancer.
AI can be the force multiplier that makes early detection accessible to all.''
\end{flushright}

%==============================================================================
% FIM DO CAPÍTULO 10
%==============================================================================
